\documentclass[../DoAn.tex]{subfiles}
\begin{document}

\begin{center}
    \Large{\textbf{TÓM TẮT NỘI DUNG BÁO CÁO}}\\
\end{center}
Theo dõi sức khỏe thai nhi trong suốt thai kỳ đóng vai trò then chốt trong phát hiện sớm các bất thường, đặc biệt là dị tật tim bẩm sinh. Trong các phương pháp không xâm lấn, điện tâm đồ thai nhi (fetal ECG) trích xuất từ tín hiệu ECG của mẹ là hướng tiếp cận an toàn và giàu thông tin; tuy nhiên đây là bài toán khó do tín hiệu mẹ--thai chồng lấn và nhiễu từ môi trường, trong khi các kỹ thuật xâm lấn tiềm ẩn rủi ro còn CTG và siêu âm Doppler vẫn hạn chế về độ chính xác lẫn độ an toàn.

Đề tài ``Ứng dụng Kolmogorov–Arnold Networks trong Transformer W-NET cho bài toán trích xuất điện tim thai nhi không xâm lấn'' đề xuất kết hợp mô hình Kolmogorov--Arnold Networks (KAN) với kiến trúc Transformer W-NET nhằm nâng cao chất lượng trích xuất tín hiệu fetal ECG từ ECG của mẹ. Bên cạnh việc đề xuất và tối ưu mô hình, đề tài còn thiết kế và triển khai một hệ thống nguyên mẫu theo thời gian thực gồm bo mạch Raspberry Pi, module thu nhận tín hiệu ADS1293 và phần mềm máy chủ FastAPI để xử lý tín hiệu và phân tích nhịp tim trực tuyến.

Kết quả thực nghiệm cho thấy mô hình đạt hiệu năng cao, đồng thời hệ thống triển khai đáp ứng tốt yêu cầu về tốc độ xử lý và khả năng giám sát trực quan, qua đó thể hiện giá trị khoa học và ý nghĩa thực tiễn trong hỗ trợ y học lâm sàng.

\begin{flushright}
Sinh viên thực hiện\\
\begin{tabular}{@{}c@{}}
\textit{(Ký và ghi rõ họ tên)}
\end{tabular}
\end{flushright}

\end{document}